\subsection{ARCH-136-PD-SECONDARY: a secondary class and its exact invariance limit}
\label{sec:arch136pdsecondary}

\paragraph{Sheet and scope.}
The divided-power structure is that of cote 136, physical p.~111,
footer 110, \texttt{ARCHIVE\_136\_scan111.tex}. The question of using
it to distinguish admissible theories concerns conditions (3)--(4)
on physical pp.~117 and 119, footers 116 and 118, and conjecture (a)
on physical p.~121, footer 120. The following identities concern
strict differential PD algebras; they do not assert a component
invariant of \(\mathcal A_{\mathbf Z}\). They use no A, M, or outside
preprint. Their input tag is none (explicit algebra).

Let \(D\) be a nonnegatively bigraded strictly commutative differential
PD algebra over \(\mathbf Z\), with differential of bidegree \((0,1)\).
An even element in its PD ideal has
\[
 d\gamma_r(x)=\gamma_{r-1}(x)\,dx.
\]
Fix a prime \(p\), an odd integer \(q\geq1\), and an exact even
element \(b=da\) of weight \(w>0\) and form degree \(q+1\).
Assume additionally that \(\gamma_p(b)=du\). Necessarily \(u\) has
weight \(pw\) and form degree \(p(q+1)-1\). Put
\[
 \Theta_p(b)=
 \bigl[p u-a\gamma_{p-1}(b)\bigr]
 \ \in\
 H^{pw,p(q+1)-1}(D)/
       pH^{pw,p(q+1)-1}(D).
 \tag{136.PD1}
\]

\paragraph{Statement.}
The class (136.PD1) is well-defined, independent of the choices of
both \(a\) and \(u\), with \(b\) fixed. It is natural under strict
weight-preserving differential PD maps. In fact these assertions
remain true with any integer \(p\geq2\) in place of a prime.

\paragraph{Proof.}
Since \(db=0\), the divided-power multiplication identity gives
\[
 d\bigl(a\gamma_{p-1}(b)\bigr)
       =b\gamma_{p-1}(b)=p\gamma_p(b).
\]
Thus the displayed representative is closed. Replacing \(u\) by
another primitive changes the class by \(p\) times a cohomology
class. If \(a\) is replaced by \(a+c\), where \(dc=0\), the change
is represented, up to sign, by \(c\gamma_{p-1}(b)\). As \(q\) is
odd,
\[
 d\bigl(ca\gamma_{p-2}(b)\bigr)
       =-(p-1)c\gamma_{p-1}(b).
\]
The cohomology class \(x=[c\gamma_{p-1}(b)]\) is therefore killed by
\(p-1\). It satisfies \(x=px\), so its image modulo \(pH\) is zero.
This establishes independence without any torsion-freeness
assumption on the groups of \(D\). A strict PD map carries all
displayed choices and identities to their counterparts in the
target, which proves naturality. \(\square\)

\paragraph{ARCH-136-PD-COEFFICIENT: a canonical vanishing.}
Let \(c\) be any closed degree-zero element of positive weight in the
PD ideal. For every \(b=da\) as above, the element
\(\gamma_p(cb)\) is exact, even without assuming \(\gamma_p(b)\)
exact, and
\[
                         \Theta_p(cb)=0.
 \tag{136.PD2}
\]
Indeed \(c^p=p!\gamma_p(c)\). The explicit primitive
\[
 u'=(p-1)!\gamma_p(c)\,a\gamma_{p-1}(b)
\]
satisfies
\[
 du'=c^p\gamma_p(b)=\gamma_p(cb),\qquad
 p u'-(ca)\gamma_{p-1}(cb)=0
\]
as cochain identities. This proves the assertion. In particular,
the operation vanishes on \(T^{[s]}b\) for every \(s>0\) in the
coefficient algebra \(S=\Gamma_{\mathbf Z}(T)\). This is an
annihilation statement about multiplying the input. It does not
identify different weight components of cohomology.

\paragraph{ARCH-136-PD-RANGE: cohomology isomorphisms do not suffice.}
Even a strict PD map inducing isomorphisms on cohomology through the
output degree and one degree beyond need not identify the full
range of (136.PD1) over its exact inputs.

\paragraph{Proof by a finite-degree example.}
Let
\[
 D=\Lambda_{\mathbf Z}(a,u)\otimes\Gamma_{\mathbf Z}(b),
 \quad |a|=1,\quad |b|=2,\quad |u|=3,
 \quad da=b,\quad db=0,\quad du=b^{[2]}.
 \tag{136.PD3}
\]
Give \(a,b\) weight one and \(u\) weight two. This is the usual free
strict PD graded algebra on its even and odd generators, with the
indicated special derivation. The square of the derivation is zero.
The element
\[
                         z=2u-ab
\]
is closed of form degree three and weight two. Let
\(E=\Lambda_{\mathbf Z}(z)\), with \(dz=0\), and map \(E\to D\) by
this displayed formula. It is a strict differential PD map.

Here is the complete low-degree calculation, with products written
in their free basis:
\[
\begin{array}{c|ccccc}
 j&0&1&2&3&4\\ \hline
 D^j&\mathbf Z&\mathbf Za&\mathbf Zb&
       \mathbf Z(ab)\oplus\mathbf Zu&
       \mathbf Zb^{[2]}\oplus\mathbf Z(au).
\end{array}
\]
The relevant differentials are
\[
 da=b,\quad d(ab)=2b^{[2]},\quad du=b^{[2]},\quad
 d(au)=bu-ab^{[2]}.
\]
The last two monomials are distinct free basis vectors in degree
five. Hence \(H^1(D)=H^2(D)=H^4(D)=0\), and
\(H^3(D)=\mathbf Zz\). Together with \(H^0(D)=\mathbf Z\), this
shows that \(E\to D\) induces isomorphisms in all form degrees
\(0\leq j\leq4\), separately in each weight.

Every exact degree-two input in \(E\) is zero because \(E^1=0\),
and \(\Theta_2(0)=0\). In \(D\), the input \(b=da\), with
primitive \(u\) of \(\gamma_2(b)\), gives
\[
                   \Theta_2(b)=[z]\neq0
                       \quad\hbox{in }H^3(D)/2H^3(D).
\]
Thus the full ranges differ. These algebras are not claimed to be
admissible simplicial theories. The example disproves the proposed
inference from cohomology isomorphisms to range invariance, not
connectedness of \(\mathcal A_{\mathbf Z}\). \(\square\)

\paragraph{Consequence for the current integral route.}
The secondary expression does record PD information beyond the
ordinary associative comparison. Its naturality with a
\emph{specified} exact input is weaker than invariance of its full
range along a strict arrow: an arbitrary nonclosed primitive in the
target need not lift. The preceding example rules out obtaining
that stronger assertion solely from the uniform cohomology
comparison. Moreover its form degree is at least \(2p-1\), so
letting \(p\) grow does not keep a fixed form degree; admissibility
allows the weight cutoff to depend on that degree. Finally
(136.PD2) prevents propagation of one fixed nonzero value to the
whole coefficient tail by multiplying the input.

A component obstruction would therefore still need an invariance
argument for actual admissible theories and a separating pair in
arbitrarily high relevant weights. Neither has been supplied here.
The construction is compatible with, but does not prove the
existence of, a completion that adjoins a primitive for
\(\gamma_p(b)\) and then a primitive for
\(p u-a\gamma_{p-1}(b)\).
